\chapter{Lane}

Winter had already settled in, and the mid-term exams had begun. The students threw themselves into work — some now taking tuitions at several places, for several subjects at once. The teachers pushed hard, reminding the class again and again that the ninth standard was a different thing altogether from the years before: if they wanted to do well in next year's board exams, the work had to start here, in the ninth itself. Shreya was in the tenth standard, and her boards fell that very year.

Shiva felt none of that pressure; he had always been good at his studies. And in any case, he had Nana at home whenever he needed help. He kept going to the coaching all the same — for the extra hours with his friends, and for the simple pleasure of being out of the house. The exams came and went, and the students finally let themselves relax. Shiva's group — the ones who went to coaching with him, along with a few others — had fallen into the habit of staying back at school a couple of hours after it let out, to play volleyball. Shiva had little skill for the game. But he loved it, and he joined them anyway, even if that mostly meant chasing down the balls that flew out of the court.

One day, playing on the court across from the school building, he caught sight of Gauri walking along the corridor. It was the first time he had seen her in days — her attendance had grown thin and irregular of late. Curiosity got the better of him, and he turned to one of his friends to ask what she was doing here, so long after school had let out.

``What can I tell you? I saw her here yesterday as well, after school got over,'' Rajveer replied.

``But why?'' Shiva pushed as if someone should know.

``She's joined a tuition here. Mr. Verma teaches Science and Maths after school. It's good for these people—they can stay in the school hostel and don't have to worry about space and all,'' one of them replied.

Vijay got in: ``Take our Robin sir, for example. He has to pay eight thousand for the room we study in.''

``And those studying there—we pay on time, but some haven't paid the fees in four or five months. I don't think you'll find another teacher like him; he never even asks for them,'' Rajveer continued.

``He earns plenty regardless. He takes our batch before going to his school, and after the school break he runs two more. He even owns a share in the school he teaches at,'' Vijay added with a smirk.

``I didn't know that,'' Rajveer exclaimed. ``He's friends with my older cousin—he told me. You lot know nothing.''

``Are you all going to keep talking, or should we start the game?'' someone called from the other side of the net.

``Yeah, start it. We're ready,'' Vijay said, quick as ever to take charge.

Shiva looked up at the second floor. It was empty. Then she stepped out of a classroom into the corridor—and she wasn't alone. Her friend Roshini was with her. Roshini stopped at the railing, seeming to take an interest in the game on the ground below. Gauri didn't. She tugged at Roshini's arm, as if asking her to come away, but Roshini held her ground and pulled her back to the railing. Shiva stood rooted. Then Gauri finally glanced down, and their eyes met. For a moment the world stopped for him—until the ball struck him in the stomach. His balance went and he dropped to the ground. Everyone burst out laughing. ``What happened to you? Can't you see a ball that big?'' Vijay wasn't pleased. But Shiva wasn't watching the game; he was watching the second floor. Then she smiled too. He broke the gaze, walked out toward the ball, threw it back in, and left the court altogether.

Vijay yelled after him, ``What happened? Don't be angry, I was only joking. Come back and play from the front if you want.''

``It's fine. I'm just tired,'' Shiva called back. ``I'll join again soon. Let me get some water first.''

Shiva walked to the tap near the building. There he saw Mr. Verma coming with his books, heading for the stairs. Shiva hurried over and greeted him. ``Good afternoon, sir.''

``Oh, hello, Shiva. What are you doing here after school?'' he asked, sounding surprised.

``I was just playing volleyball with my friends,'' he answered.

``Well, congratulations—you got the highest marks in Maths in the mid-term exams,'' Mr. Verma said, pride and warmth in his voice.

``Thank you, sir. It's all because of you,'' he responded, head bowed.

``All right, carry on with your friends—I have a class to take,'' Mr. Verma said, and turned to go up the stairs.

He had climbed only a few steps when Shiva hurried after him and called out, ``Sir.''

``Yes?'' Mr. Verma stopped and turned back. Shiva climbed the stairs up to him. ``Sir, could you please teach me Science? I don't think I did well in the exams. I feel my marks will slip if things go on the way they are.''

``No, no. How could you need a tuition? You've always been top of the class in nearly every subject.''

``Yes, sir. But ever since ninth standard began, Science has changed completely. Now it's split into Physics, Chemistry, and Biology instead of just one subject. I think I need a little help.''

``Sure, why not. Join us after class. I'm already teaching Roshini and Gauri—both Science and Maths. You could even help them with Maths while you learn Science.''

``Yes, sir. I'll be ready tomorrow,'' he said. His face broke into a deep smile. His whole world had turned upside down.

He headed back to the match, glancing once at the second floor—empty again now—though the thought of tomorrow was already enough to lift him past the sting of the fall. ``Are you better now?'' Vijay asked, concern in his eyes as he walked over. ``Yeah, perfect now,'' Shiva replied, with such enthusiasm that Vijay smiled too. 